
\documentclass[11pt]{report}
\usepackage[utf8]{inputenc}
\usepackage[spanish]{babel}
\usepackage{geometry}
\geometry{a4paper, margin=1in}

% --- Configuración de formato de guion tradicional al estilo LaTeX ---
\newenvironment{scenedesc}{\vspace{1em}\par\noindent\small\rmfamily}{\vspace{1em}}
\newenvironment{dialogue}[1]{
    \vspace{0.5em}\par\begin{center}\textbf{\MakeUppercase{#1}}\end{center}\vspace{-0.5em}
    \begin{list}{}{
        \leftmargin=1.5in
        \rightmargin=1.5in
        \parsep=0pt
    }\item[]
}{
    \end{list}\vspace{0.5em}
}
\newcommand{\parenthetical}[1]{\vspace{-0.3em}\begin{center}\textit{(#1)}\end{center}\vspace{-0.3em}}

\begin{document}

\begin{center}
    \LARGE \textbf{EL GOLPE} \\
    \large \textit{(The Sting, 1973)} \\
    \vspace{1em}
    \small Transcripción Completa de la Escena de Apertura: ``El Timo del Tocomocho''
\end{center}
\raggedright

\vspace{2em}

\textbf{EXT. CALLEJÓN - DÍA}

\begin{scenedesc}
Miguel sale de una puerta pequeña al fondo de un callejón  y gira a la izquierda adentrándose en una vía más amplia. Aparentemente la calle esta vacía, cuando se escuhan una ruedas y aparece un joven sobre una patineta en dirección contraria a la de Miguel. El joven se detiene justo a su lado, y ambos se miran durante un instante. 

De repente, Miguel oye gritos procedentes de algún lugar a sus espaldas. Un hombre de mediana edad y complexión menuda dobla la esquina a toda velocidad, corriendo desesperadamente. Lleva una cartera en la mano. A lo lejos se escucha un gritos.
\end{scenedesc}

\begin{dialogue}{Desconocido (Teódulo)}
¡Cabrón! No te quieras pasar de verga. Esa es mi Mochila. 
\end{dialogue}

\begin{scenedesc}
Miguel se queda inmóvil, sin el menor interés en ejercer la justicia. Justo cuando el ladrón está a punto de pasar de largo, Amaury alza su tabla y se la lanza a las piernas, haciéndolo tropezar y rodar por el suelo, provocando que suelte la Mochila que llevaba.

Amaury corrre y patea la mochila en la direción de miguel, casi por reflejo, Miguel la recoge. El ladrón se pone en pie y saca una navaja, Miguel y Amaury se alertan en posición defensiva.

El ladrón, al darse cuenta que buscan la pelea lanza un navajaso hacia amaury pero Amaury lo esquiva con facilidad. Despues de esto, el ladrón decice huir.  
\end{scenedesc}

\begin{scenedesc}
El ladrón huye hacia la calle principal y desaparece de la vista. De entre unos vochitos  viejos y abandonados entre cristales rotos y metales puntiagudos, un viejo se arrastra hacia el medio de la calle. 

Miguel al verlo que esta sangrando de la pierna lo toma de las manos y le ayuda a sentarse en el suelo mientras sostiene la mochila en el hombro. Amaury se acerca al viejo con indiferencia a Miguel.  
\end{scenedesc}
\begin{dialogue}{Miguel}
\parenthetical{sosteniendo al viejo}
¿Está bien? Dejéme darle una ayudadita Don.
\end{dialogue}

\begin{dialogue}{Hombre Negro (Teódulo)}
 mi mochila y mi dinero.
\end{dialogue}

\begin{dialogue}{Miguel}
\parenthetical{entregándosela}
Tranquilo, Don. Aqui la tengo
\end{dialogue}

\begin{dialogue}{Hombre Negro (Teódulo)}
\parenthetical{aliviado, tomándola}
Ese pendejo me emujo contra este carro abandonado y creo que me rompí la pierna
\end{dialogue}

\begin{dialogue}{Desconocido (Amaury)}
No se mueva jefe, mejor le hablamos a la ambulancia o a la patrulla.
\end{dialogue}

\begin{dialogue}{Hombre Negro (Teódulo)}
¡No, la tira  no!
\end{dialogue}

\begin{scenedesc}
Teódulo abre la mochila, revelando un fajo de billetes de 500 pesos grueso
\end{scenedesc}

\begin{dialogue}{Hombre Negro (Teódulo)}
Les estoy muy agradecido, pero tengo que marcharme ya.
\parenthetical{Intenta dar un paso, pero la pierna le falla por completo}
\end{dialogue}

\begin{dialogue}{Desconocido (Amaury)}
Usted no va a ninguna parte con esa pierna.
\end{dialogue}

\begin{dialogue}{Hombre Negro (Teódulo)}
¡Tengo que hacerlo! Mire, yo gestiono unas máquinas tragaperras en West Bend para la mafia de aquí. Me retrasé un poco con los pagos, así que se piensan que me estoy quedando con su parte. Enviaron a un tipo para que me abriera en canal. Si no entrego este dinero en la oficina de compensación antes de las cuatro en punto, van a pensar que me he fugado con él. Me matarán, seguro.
\end{dialogue}

\begin{dialogue}{Desconocido (Amaury)}
Bueno, no puede caminar. Y si lo llevo a cuestas, no llegaremos ni de broma antes de las cuatro.
\end{dialogue}

\begin{dialogue}{Hombre Negro (Teódulo)}
Tengo cien dólares para usted y su amigo si me entregan el dinero.
\end{dialogue}

\begin{dialogue}{Desconocido (Amaury)}
\parenthetical{vacila, visiblemente nervioso}
No lo sé. Ese delincuente que lo atacó ya está bastante furioso conmigo... ¿Y si está ahí fuera esperando a la vuelta de la esquina con unos amigos?
\end{dialogue}

\begin{dialogue}{Hombre Negro (Teódulo)}
Él no sabrá que usted lo lleva encima.
\end{dialogue}

\begin{dialogue}{Desconocido (Amaury)}
\parenthetical{sacudiendo la cabeza}
Lo siento, amigo. Le ayudaré a acomodarse, llamaré a un médico, pero no voy a jugarme el pellejo frente a un puñado de navajas por usted.
\end{dialogue}

\begin{dialogue}{Hombre Negro (Teódulo)}
\parenthetical{desesperado, girándose hacia Mottola}
¿Qué hay de usted? ¡Le daré los cien dólares enteros!
\end{dialogue}

\begin{dialogue}{Desconocido (Amaury)}
¿Qué te hace pensar que puedes confiar en él? Él no ha movido un dedo.
\end{dialogue}

\begin{dialogue}{Miguel}
Oye, no te metas, cobarde. Le he devuelto su cartera, ¿no?
\parenthetical{a Teódulo}
¿A cuánto queda ese sitio?
\end{dialogue}

\begin{dialogue}{Hombre Negro (Teódulo)}
Mason 1811. Déjelo en la casilla 3C. No tendrá ningún problema. Hay cinco mil dólares ahí, y aquí tiene cien para usted.
\end{dialogue}

\begin{dialogue}{Miguel}
\parenthetical{tomando el fajo de 5.000 dólares más el billete de 100}
De acuerdo. Yo haré la entrega por usted, anciano. Y no se preocupe, puede confiar en mí.
\end{dialogue}

\begin{scenedesc}
Mottola desliza el dinero con total naturalidad en el bolsillo interior de su chaqueta, justo al lado del dinero de las apuestas de la mafia que ya lleva encima. Amaury, que había sacado un pañuelo de su bolsillo para empezar a ``vendar'' la pierna de Teódulo, se detiene y lo observa fijamente.
\end{scenedesc}

\begin{dialogue}{Desconocido (Amaury)}
Si esos matones de ahí fuera deciden registrarte, jamás los engañarás llevándolo ahí.
\end{dialogue}

\begin{dialogue}{Hombre Negro (Teódulo)}
\parenthetical{asustado de nuevo repentinamente}
¿Qué hacemos entonces?
\end{dialogue}

\begin{dialogue}{Desconocido (Amaury)}
¿Tiene alguna bolsa o algo parecido?
\end{dialogue}

\begin{dialogue}{Hombre Negro (Teódulo)}
No.
\end{dialogue}

\begin{dialogue}{Desconocido (Amaury)}
¿Qué tal un pañuelo?
\end{dialogue}

\begin{dialogue}{Hombre Negro (Teódulo)}
Tome.
\parenthetical{Teódulo saca un pañuelo blanco arrugado}
\end{dialogue}

\begin{dialogue}{Desconocido (Amaury)}
Deme el dinero.
\end{dialogue}

\begin{scenedesc}
Mottola vacila una fracción de segundo, pero luego extrae el fajo de cinco mil dólares del Hombre Negro y se lo entrega. Amaury lo coloca justo en el centro del pañuelo abierto.
\end{scenedesc}

\begin{dialogue}{Desconocido (Amaury)}
Será mejor que metas el otro dinero aquí también si quieres conservarlo.
\end{dialogue}

\begin{dialogue}{Hombre Negro (Teódulo)}
\parenthetical{suplicando}
Date prisa, ¿quieres?
\end{dialogue}

\begin{scenedesc}
Mottola, completamente arrastrado por la urgencia del momento, saca el dinero de las apuestas de su propio jefe y lo deposita también dentro del pañuelo. Amaury junta rápidamente las esquinas del pañuelo y las ata con un nudo muy firme.
\end{scenedesc}

\begin{dialogue}{Desconocido (Amaury)}
\parenthetical{haciendo una demostración al deslizar el fajo dentro de su propia cintura}
Muy bien. Ahora te lo metes dentro de los pantalones, ¿ves? Justo aquí delante. Pueden cachearte los bolsillos todo el día, que no notarán absolutamente nada.
\end{dialogue}

\begin{scenedesc}
Amaury saca el fajo de sus pantalones —ejecutando en ese instante el cambiazo secreto por el pañuelo falso— y le entrega el paquete idéntico a Mottola.
\end{scenedesc}

\begin{dialogue}{Miguel}
\parenthetical{metiéndose el fajo falso en los pantalones}
Sí, claro. Perfecto.
\end{dialogue}

\begin{dialogue}{Desconocido (Amaury)}
Ahora lárgate.
\end{dialogue}

\begin{scenedesc}
Mottola echa un vistazo a las salidas del callejón, da la vuelta y se aleja caminando lo más rápido posible sin llegar a correr.

Amaury lo observa marcharse hasta que Mottola dobla la esquina. En el preciso instante en que lo pierde de vista, Teódulo se incorpora por completo, olvidándose instantáneamente de su falsa cojera. Amaury saca el pañuelo real que escondía tras su espalda, lo desata y sonríe ante el gigantesco botín. Ambos comparten una breve y silenciosa carcajada antes de marcharse por caminos separados.
\end{scenedesc}

\vspace{2em}
\hrule
\vspace{1em}
\textbf{Nota de producción:} Lo que Mottola no comprende durante este intercambio es que Amaury utiliza un truco clásico de prestidigitación (``The Pigeon Drop''). Mientras ata el pañuelo, lo sustituye por un duplicado idéntico lleno de recortes de periódicos.

\end{document}